Musical Information Retrieval (MIR) has emerged as an important research area in machine learning, with automatic note segmentation and recognition remaining challenging tasks. While considerable progress has been made for Western instruments such as the piano, comparatively little research has focused on Indian classical music and its instruments, resulting in limited datasets and methodologies for automatic swara transcription. This study proposes an end-to-end pipeline for automatic segmentation and swara classification of monophonic Indian classical flute melodies.

The proposed framework consists of two stages: note segmentation and swara classification. A dataset of 748 individual flute notes was collected through a combination of self-recordings and publicly available recordings. These notes were then used to generate a synthetic melody dataset that closely resembles continuous flute performances. An RCNN with an attention mechanism was employed to identify note boundaries within continuous melodies, while a CNN with an attention mechanism was trained to classify segmented notes using Constant-Q Transform (CQT) features.

To investigate the effectiveness of different acoustic representations, multiple feature extraction techniques, including MFCC, Log-Mel Spectrogram, CQT, pitch, RMS energy, and their combinations, were evaluated for the segmentation stage. Experimental results demonstrate that CQT-based features achieved the best overall performance. The proposed segmentation model achieved a segmentation accuracy of 92.86\%, while the classification model achieved 97.80\% swara classification accuracy on correctly segmented melodies.

The results demonstrate that accurate end-to-end swara transcription can be achieved using a relatively small dataset when combined with synthetic data generation, careful model optimization, and appropriate feature representation. The proposed framework provides a practical and scalable approach for automatic transcription of monophonic Indian classical flute melodies and has the potential to be adapted to other monophonic musical instruments with minimal annotated data.
